\subsection{混合预条件子}

\begin{figure}[H]
  \centering
  \begin{subfigure}[b]{\textwidth}
    \centering
    \includegraphics[width=0.8\textwidth]{sheet/pdf/single_chain_preconditioner_C.pdf}
    \caption{平均预计算时间随线段数的变化}
    \label{fig:pre}
  \end{subfigure}

  \begin{subfigure}[b]{\textwidth}
    \centering
    \includegraphics[width=0.8\textwidth]{sheet/pdf/single_chain_preconditioner_A.pdf}
    \caption{平均 PCG 迭代次数随线段数的变化}
  \end{subfigure}

  \begin{subfigure}[b]{\textwidth}
    \centering
    \includegraphics[width=0.8\textwidth]{sheet/pdf/single_chain_preconditioner_B.pdf}
    \caption{平均求解时间随线段数的变化}
    \label{fig:compute}
  \end{subfigure}

  \caption[本文预条件子与其他预条件子比较]{本文预条件子与其他预条件子比较。
  在重力作用下的螺旋线场景中，对不同预条件子运行 100 个时间步，并在每个时间步最多进行 5 次牛顿迭代。
  本文预条件子的构造开销略高（a）；但在线性系统求解阶段，PCG 迭代次数（b）与总耗时（c）均明显小于对角或块对角预条件子。随着线段数增加，该优势更为明显。}
  \label{fig:compare_preconditioner}
  \vspace{-0.12in}
\end{figure}

合适的预条件子对整体性能至关重要。常见做法是直接取 $\mathbf{A}$ 的对角或块对角部分构造预条件子，因为这种方式实现简单、构造成本低。例如，RedMax~\cite{wang_redmax_2019} 就采用了块对角策略。然而，本文发现这类常见近似在当前问题上并不总是理想，下面构造的预条件子通常能获得更好的收敛表现。

具体来说，链式结构中的 $\mathbf{A}$ 可以表示为：
\begin{equation}
  \begin{aligned}
    \mathbf{A} & = \J^{\top} \mathbf{H} \J =
    \begin{bmatrix}
      \J_{\mathbf{p, s}}^{\top} &
      \J_{\mathbf{q, s}}^{\top}
    \end{bmatrix}
    \begin{bmatrix}
      \mathbf{H}_{\mathbf{p}} & \mathbf{0}              \\
      \mathbf{0}              & \mathbf{H}_{\mathbf{q}}
    \end{bmatrix}
    \begin{bmatrix}
      \J_{\mathbf{p, s}} \\
      \J_{\mathbf{q, s}}
    \end{bmatrix}                                                          \\
               & =\J^{\top}_{\mathbf{p, s}} \mathbf{H}_{\mathbf{p}} \J_{\mathbf{p, s}} +
    \J^{\top}_{\mathbf{q, s}} \mathbf{H}_{\mathbf{q}} \J_{\mathbf{q, s}}.
  \end{aligned}
  \label{eq:A}
\end{equation}

进一步分析可得，$\mathbf{A}$ 具有如下结构：
\begin{equation}
  \begin{aligned}
    \J_{\mathbf{q, s}}^{\top} \mathbf{H}_{\mathbf{q}} \J_{\mathbf{q, s}} & =
    \begin{bmatrix}
      \mathbf{0} & \mathbf{0}                                                                   \\
      \mathbf{0} & \J_{\mathbf{q, \omega}}^{\top}\mathbf{H}_{\mathbf{q}}\J_{\mathbf{q, \omega}}
    \end{bmatrix} =
    \begin{bmatrix}
      \mathbf{0} & \mathbf{0} \\
      \mathbf{0} & \TriDiag
    \end{bmatrix}
    ,                                                                        \\
    \J^{\top}_{\mathbf{p, s}} \mathbf{H}_{\mathbf{p}} \J_{\mathbf{p, s}} & =
    \mathrm{Diag}\{\J^{i}_{\mathbf{p,s}}\}^{\top} \cdot
    \underbrace{
    \UpperTri \cdot \mathbf{H}_{\mathbf{p}}\cdot \LowerTri}_{\bar{\mathbf{H}}_{\mathbf{p}}}
    \cdot \mathrm{Diag}\{\J^{i}_{\mathbf{p,s}}\},
  \end{aligned}
  \label{eq:JTHJ}
\end{equation}
其中 $\J_{\mathbf{q, s}} =
  \begin{bmatrix}
    \mathbf{0}, \J_{\mathbf{q,\omega}}
  \end{bmatrix}
$且 $\J_{\mathbf{p, s}} =
  \begin{bmatrix}
    \I, \J_{\mathbf{p, \omega}}
  \end{bmatrix}$。

\cref{eq:JTHJ} 表明，$ \J_{\q,\omega}^\top \mathbf{H}_\q \J_{\q, \omega}$ 具有块三对角结构。受此启发，本文提取 $\J_{\mathbf{q, s}}^{\top} \mathbf{H}_{\mathbf{q}} \J_{\mathbf{q, s}}$ 的完整块三对角部分，以及 $\J^{\top}_{\mathbf{p, s}} \mathbf{H}_{\mathbf{p}} \J_{\mathbf{p, s}}$ 的对角线部分来构造混合预条件子 $\mathbf P$，以尽可能保留变换 $\mathcal T$ 后的信息：
\begin{equation}
  \setlength{\arraycolsep}{0.5pt}
  \mathbf{P} =\left(\operatorname{blk-diag}\left(\J_{\p, \s}^\top \mathbf{H}_{\p} \J_{\p, \s}\right) + \begin{bmatrix} \mathbf{0} & \mathbf{0} \\ \mathbf{0}  & \J_{\q,\omega}^\top \mathbf{H}_\q \J_{\q, \omega} \end{bmatrix}\right).
  \label{eq:preconditioner_structure}
\end{equation}

这种构造并非任意截断系统矩阵，而是有针对性地保留当前问题中最关键的结构成分。对于与旋转变量相关的 $\mathbf q$ 部分，块三对角结构反映了沿链条传播的局部弯扭耦合，因此本文完整保留该结构。相比之下，$\mathbf p$ 部分在链式表示下更多体现为由质量和惯性诱导的平移响应；在预条件子中仅保留其对角部分，意味着保留局部主导尺度，同时忽略较弱的远距离耦合。这样既避免预条件子退化为近似原矩阵的稠密对象，又能在力学上抓住最主要的刚性来源。

换言之，混合预条件子的设计逻辑是：对决定高频刚性传播的块结构尽量完整保留，对只提供局部尺度信息、但不主导收敛行为的部分采用更轻量的近似。这种取舍围绕链式杆系统的实际耦合模式展开，使预条件子能够在不显著增加应用代价的前提下，同时兼顾正定性、结构稀疏性和收敛效果。

该混合预条件子具有较好的性质：它是对称正定矩阵，且构造与应用都可在线性时间 $\mathcal{O}(n)$ 内完成。

\inputbody{paper1/app_invertible_proof}

\textbf{$\mathcal{O}(n)$ 时间构造：} 考虑到这种特殊的 [块对角矩阵]-[块矩阵]-[块对角矩阵] 乘法：
\begin{equation}
  \label{eq:bmb_multiplication}
  \setlength{\arraycolsep}{1.2pt}
  \begin{bmatrix}
    \mathbf{A}_{1} &                &        &                \\
                   & \mathbf{A}_{2} &        &                \\
                   &                & \ddots &                \\
                   &                &        & \mathbf{A}_{n}
  \end{bmatrix}
  \begin{bmatrix}
    \mathbf{B}_{11} & \mathbf{B}_{12} & \cdots & \mathbf{B}_{1n} \\
    \mathbf{B}_{21} & \mathbf{B}_{22} & \cdots & \mathbf{B}_{2n} \\
    \vdots          & \vdots          & \ddots & \vdots          \\
    \mathbf{B}_{n1} & \mathbf{B}_{n2} & \cdots & \mathbf{B}_{nn} \\
  \end{bmatrix}
  \begin{bmatrix}
    \mathbf{C}_{1} &                &        &                \\
                   & \mathbf{C}_{2} &        &                \\
                   &                & \ddots &                \\
                   &                &        & \mathbf{C}_{n}
  \end{bmatrix}.
\end{equation}
结果矩阵 $\mathbf A \mathbf B \mathbf C$ 的 $(i, j)$ 块为 $\mathbf{A}_{i} \mathbf{B}_{ij} \mathbf{C}_{j}$。利用这一性质，当矩阵 $\mathbf A$、$\mathbf B$ 和 $\mathbf C$ 已知时，可在线性时间内构造其块三对角部分或块对角部分。

回到本文预条件子，$\J_{\q,\omega}^\top \mathbf{H}_\q \J_{\q, \omega}$ 满足这种结构，因为 $\J_{\mathbf{q, \omega}}$ 是块对角矩阵 (\cref{eq:Jqs_pattern})。利用~\cref{eq:JTHJ} 中的乘法结合律，$\J_{\p,\s}^\top \mathbf{H}_\p \J_{\p, \s}$ 也可转换为~\cref{eq:bmb_multiplication} 的结构。因此，问题转化为如何在线性时间内得到 $\bar{\mathbf{H}}_{\mathbf{p}}$。注意到 $\mathbf{H}_{\mathbf{p}}$ 是集中块对角矩阵：
\begin{equation}
  \bar{\mathbf{H}}_{\mathbf{p}} = \UpperTri \mathbf{H}_{\mathbf{p}} \LowerTri =
  \begin{bmatrix}
    \sum_{i=1}^{n} \mathbf{h}_{i} & \sum_{i=2}^{n} \mathbf{h}_{i} & \sum_{i=3}^{n} \mathbf{h}_{i} & \cdots & \mathbf{h}_{n} \\
    \sum_{i=2}^{n} \mathbf{h}_{i} & \sum_{i=2}^{n} \mathbf{h}_{i} & \sum_{i=3}^{n} \mathbf{h}_{i} & \cdots & \mathbf{h}_{n} \\
    \sum_{i=3}^{n} \mathbf{h}_{i} & \sum_{i=3}^{n} \mathbf{h}_{i} & \sum_{i=3}^{n} \mathbf{h}_{i} & \cdots & \mathbf{h}_{n} \\
    \vdots                        & \vdots                        & \vdots                        & \ddots & \vdots         \\
    \mathbf{h}_{n}                & \mathbf{h}_{n}                & \mathbf{h}_{n}                & \cdots & \mathbf{h}_{n} \\
  \end{bmatrix}.
\end{equation}
其中 $\mathbf{h}_{i}$ 是 $\mathbf{H}_{\mathbf{p}}$ 的第 $i$ 个对角块。因此，$\bar{\mathbf{H}}_{\mathbf{p}}$ 仅包含 $n$ 个不同的非零块，可采用类似前缀和的串行递推方式在线性时间内计算。总的来说，该预条件子只含有 $\mathcal{O}(n)$ 个非零块，因此可以逐块显式构造。

\textbf{$\mathcal{O}(n)$ 时间求解：} 由于预条件子具有标准的块三对角结构，使用 Thomas 算法~\cite{thomas1949elliptic} 在 $\mathcal{O}(n)$ 时间内求解。

虽然本文预条件子的构造过程比对角或块对角预条件子更复杂，但这种有针对性的设计通常能够在求解速度上带来数量级的提升（见图~\ref{fig:compare_preconditioner}）。
本文还将 PCG 方法与通过矩阵分解直接构造 $\mathbf A^{-1}$ 的方式进行了比较。由于 $\mathbf A$ 具有稠密结构，矩阵分解阶段（图~\ref{fig:pre}）和矩阵--向量乘法阶段（图~\ref{fig:compute}）都会消耗更多时间。

{\color{black}
此外，在实验设置下，PCG 求解器呈现出近似线性的复杂度。在线性系统求解阶段，平均耗时随线段数增加而近似线性增长（见图~\ref{fig:linear_increasing}），这说明本文方法对段数具有较好的可扩展性。需要注意的是，这种线性复杂度并不直接适用于整个时间步，因为~\cref{eq:collision_constraint} 中的碰撞约束在本文的简化表示下仍是非线性的，不同帧所需的牛顿迭代次数也可能不同。
}

\begin{figure}[H]
  \centering
  \includegraphics[width=\textwidth]{sheet/pdf/linear_increasing.pdf}
  \caption[平均时间成本]{在图~\ref{fig:visual_damping_penalty} 中的重力作用下模拟螺旋线。PCG 求解线性系统使用几乎固定的迭代次数。每步的时间成本关于段数几乎是线性的。}
  \label{fig:linear_increasing}
\end{figure}

根据~\cref{eq:A}，$\mathbf{A}$ 由两部分相加得到。随着杨氏模量增大，$\J_{\q,\mathbf{s}}^\top\mathbf{H}_{\mathbf q} \J_{\q, \mathbf{s}}$ 在其中会占据越来越主导的地位。因此，在这种情况下，保留完整 $\J_{\q,\mathbf{s}}^\top\mathbf{H}_{\mathbf q} \J_{\q, \mathbf{s}}$ 的预条件子能够更好地逼近 $\mathbf A$，从而获得更好的收敛率（如图~\ref{A} 所示）。相比之下，简单的块对角预条件子丢弃了该项中的许多大非零块，因此收敛率较差（如图~\ref{B} 所示）。

这也解释了图~\ref{fig:compare_preconditioner} 中看似矛盾的一点：本文预条件子的预计算时间虽然更高，但总耗时反而更低。其更主要的原因不在于实现层面的技巧差异，而在于预条件子是否保留了对迭代收敛影响更大的矩阵结构信息。若预处理仅停留在对角或块对角层面，那么每次应用虽然便宜，但 PCG 仍需经历更多迭代才能逐步识别并修正链式系统中的主导耦合；相反，本文预条件子在预处理阶段就已经把最关键的块三对角弯扭结构纳入近似之中，因此能够以更少的迭代数弥补构造开销。

同样地，图~\ref{fig:PCG_iter_Young} 中随杨氏模量增大而不断扩大的性能差距，也并非偶然现象。随着材料变硬，系统中由弯扭刚度诱导的主导项会越来越集中在 $\J_{\q,\mathbf{s}}^\top\mathbf{H}_{\mathbf q} \J_{\q, \mathbf{s}}$ 这部分结构里。此时，是否保留这一部分的完整块三对角信息，会直接决定预条件子对原矩阵的逼近质量。也就是说，材料越硬，本文所保留的结构就越接近真实的主导刚性模式，因此其收敛优势会表现得更加明显。

\begin{figure}[H]
  \centering
  \vspace{1.2cm}
  \begin{subfigure}[b]{0.48\textwidth}
    \centering
    \includegraphics[width=\textwidth]{sheet/pdf/pcg_mu_our.pdf}
    \caption{本文预条件子}
    \label{A}
  \end{subfigure}
  \centering
  \begin{subfigure}[b]{0.48\textwidth}
    \centering
    \includegraphics[width=\textwidth]{sheet/pdf/pcg_mu_blk.pdf}
    \caption{块对角预条件子}
    \label{B}
  \end{subfigure}

  \caption[PCG 迭代次数与杨氏模量的关系]{随着杨氏模量的增加，本文预条件子的优势会更加明显。}
  \label{fig:PCG_iter_Young}
\end{figure}
